\chapter{Architecture of ROS2 Systems}

\begin{lernziele}
After this chapter, you can design a simple ROS2 architecture for the autonomous robot. You understand node granularity, communication structure, failure cases, parameters, and lifecycle considerations.
\end{lernziele}

\section{Introduction}
A ROS2 architecture describes which nodes a system consists of and how these nodes communicate with each other. It is the technical counterpart to many SysML structure and interface decisions.

\section{Typical Node Structure}
A reasonable first architecture consists of the following nodes:

\begin{itemize}
    \item \texttt{camera\_node}
    \item \texttt{lidar\_node}
    \item \texttt{slam\_node}
    \item \texttt{navigation\_node}
    \item \texttt{object\_detector\_node}
    \item \texttt{battery\_monitor\_node}
    \item \texttt{motion\_control\_node}
    \item \texttt{user\_interface\_node}
    \item \texttt{robot\_state\_node}
\end{itemize}

\section{Communication Structure}
\begin{lstlisting}[style=plaintext]
camera_node -> /camera/image_raw -> object_detector_node
lidar_node  -> /scan             -> slam_node
slam_node   -> /map              -> navigation_node
navigation_node -> /cmd_vel      -> motion_control_node
battery_monitor_node -> /battery_state -> navigation_node
object_detector_node -> /detected_objects -> navigation_node
robot_state_node -> /robot/status -> user_interface_node
\end{lstlisting}

Figure~\ref{fig:kap15-node-graph} shows the complete node graph with all nine nodes and their topic connections.

\begin{figure}[H]
  \centering
  \adjustbox{max width=\linewidth,max totalheight=0.6\textheight}{% =====================================================================
% Figure: ROS2 node graph of the autonomous indoor robot
% Matches Chapter 15, section "Communication Structure"
%
% Included in the book, e.g. in chapters_en/15_ros2_architektur.tex:
%
%   \begin{figure}[H]
%     \centering
%     \resizebox{\textwidth}{!}{% =====================================================================
% Figure: ROS2 node graph of the autonomous indoor robot
% Matches Chapter 15, section "Communication Structure"
%
% Included in the book, e.g. in chapters_en/15_ros2_architektur.tex:
%
%   \begin{figure}[H]
%     \centering
%     \resizebox{\textwidth}{!}{% =====================================================================
% Figure: ROS2 node graph of the autonomous indoor robot
% Matches Chapter 15, section "Communication Structure"
%
% Included in the book, e.g. in chapters_en/15_ros2_architektur.tex:
%
%   \begin{figure}[H]
%     \centering
%     \resizebox{\textwidth}{!}{\input{figures_en/fig_kap15_node_graph}}
%     \caption{Node graph of the autonomous indoor robot, with the key
%       topics between the sensor, processing, navigation, and actuator
%       nodes.}
%     \label{fig:kap15-node-graph}
%   \end{figure}
%
% Required packages (already loaded in the book via main_en.tex): tikz, graphicx
% Required TikZ libraries: positioning, arrows.meta, shadows, calc
% =====================================================================

\input{figures_en/fig_kap15_node_graph.tikz}
}
%     \caption{Node graph of the autonomous indoor robot, with the key
%       topics between the sensor, processing, navigation, and actuator
%       nodes.}
%     \label{fig:kap15-node-graph}
%   \end{figure}
%
% Required packages (already loaded in the book via main_en.tex): tikz, graphicx
% Required TikZ libraries: positioning, arrows.meta, shadows, calc
% =====================================================================

\begin{tikzpicture}[
    scale=0.74, transform shape,
    every node/.style={font=\small},
    node distance=10mm and 14mm,
    ros/.style={
        rectangle, rounded corners=2pt, draw=orange!60!black, thick,
        minimum width=32mm, minimum height=10mm,
        align=center, fill=orange!15, drop shadow
    },
    flow/.style={-{Stealth[length=2.2mm]}, thick, orange!70!black},
    lbl/.style={font=\footnotesize, fill=white, inner sep=1pt}
]

% --- Column 1: sensor nodes (left) ---
\node[ros] (camera)  at (0, 6)  {\texttt{camera\_node}};
\node[ros] (lidar)   at (0, 3)  {\texttt{lidar\_node}};
\node[ros] (battery) at (0, 0)  {\texttt{battery\_monitor\_}\\\texttt{node}};

% --- Column 2: processing ---
\node[ros] (detector) at (5.6, 6) {\texttt{object\_detector\_}\\\texttt{node}};
\node[ros] (slam)     at (5.6, 3) {\texttt{slam\_node}};

% --- Column 3: central ---
\node[ros] (nav) at (10.2, 3.5) {\texttt{navigation\_node}};

% --- Column 4: right ---
\node[ros] (motion) at (14.8, 5)   {\texttt{motion\_control\_}\\\texttt{node}};
\node[ros] (ui)      at (14.8, 0.5) {\texttt{user\_interface\_}\\\texttt{node}};

% --- additional node (status) ---
\node[ros] (state) at (10.2, -1.5) {\texttt{robot\_state\_node}};

% --- Edges ---
\draw[flow] (camera.east)  -- node[lbl, midway, above] {/camera/image\_raw} (detector.west);
\draw[flow] (lidar.east)   -- node[lbl, midway, above] {/scan}              (slam.west);
\draw[flow] (slam.east)    -- node[lbl, midway, above] {/map}               (nav.north west);
\draw[flow] (nav.east)     -- node[lbl, midway, above] {/cmd\_vel}          (motion.west);
\draw[flow] (battery.east) .. controls (5,0) and (7,2) .. node[lbl, pos=0.55, below] {/battery\_state} (nav.south west);
\draw[flow] (detector.east) -- node[lbl, midway, above] {/detected\_objects} (nav.north east);
\draw[flow] (state.east) -- node[lbl, midway, above] {/robot/status} (ui.south west);

% --- Legend ---
\node[ros, minimum width=8mm, minimum height=5mm, font=\small] (leg1) at (3.2,-3.4) {};
\node[right=2mm of leg1, font=\small, align=left] {ROS2 node};
\draw[flow] (8.4,-3.4) -- ++(12mm,0);
\node[font=\small, align=left, anchor=west] at (9.8,-3.4) {Topic (message flow)};

\end{tikzpicture}

}
%     \caption{Node graph of the autonomous indoor robot, with the key
%       topics between the sensor, processing, navigation, and actuator
%       nodes.}
%     \label{fig:kap15-node-graph}
%   \end{figure}
%
% Required packages (already loaded in the book via main_en.tex): tikz, graphicx
% Required TikZ libraries: positioning, arrows.meta, shadows, calc
% =====================================================================

\begin{tikzpicture}[
    scale=0.74, transform shape,
    every node/.style={font=\small},
    node distance=10mm and 14mm,
    ros/.style={
        rectangle, rounded corners=2pt, draw=orange!60!black, thick,
        minimum width=32mm, minimum height=10mm,
        align=center, fill=orange!15, drop shadow
    },
    flow/.style={-{Stealth[length=2.2mm]}, thick, orange!70!black},
    lbl/.style={font=\footnotesize, fill=white, inner sep=1pt}
]

% --- Column 1: sensor nodes (left) ---
\node[ros] (camera)  at (0, 6)  {\texttt{camera\_node}};
\node[ros] (lidar)   at (0, 3)  {\texttt{lidar\_node}};
\node[ros] (battery) at (0, 0)  {\texttt{battery\_monitor\_}\\\texttt{node}};

% --- Column 2: processing ---
\node[ros] (detector) at (5.6, 6) {\texttt{object\_detector\_}\\\texttt{node}};
\node[ros] (slam)     at (5.6, 3) {\texttt{slam\_node}};

% --- Column 3: central ---
\node[ros] (nav) at (10.2, 3.5) {\texttt{navigation\_node}};

% --- Column 4: right ---
\node[ros] (motion) at (14.8, 5)   {\texttt{motion\_control\_}\\\texttt{node}};
\node[ros] (ui)      at (14.8, 0.5) {\texttt{user\_interface\_}\\\texttt{node}};

% --- additional node (status) ---
\node[ros] (state) at (10.2, -1.5) {\texttt{robot\_state\_node}};

% --- Edges ---
\draw[flow] (camera.east)  -- node[lbl, midway, above] {/camera/image\_raw} (detector.west);
\draw[flow] (lidar.east)   -- node[lbl, midway, above] {/scan}              (slam.west);
\draw[flow] (slam.east)    -- node[lbl, midway, above] {/map}               (nav.north west);
\draw[flow] (nav.east)     -- node[lbl, midway, above] {/cmd\_vel}          (motion.west);
\draw[flow] (battery.east) .. controls (5,0) and (7,2) .. node[lbl, pos=0.55, below] {/battery\_state} (nav.south west);
\draw[flow] (detector.east) -- node[lbl, midway, above] {/detected\_objects} (nav.north east);
\draw[flow] (state.east) -- node[lbl, midway, above] {/robot/status} (ui.south west);

% --- Legend ---
\node[ros, minimum width=8mm, minimum height=5mm, font=\small] (leg1) at (3.2,-3.4) {};
\node[right=2mm of leg1, font=\small, align=left] {ROS2 node};
\draw[flow] (8.4,-3.4) -- ++(12mm,0);
\node[font=\small, align=left, anchor=west] at (9.8,-3.4) {Topic (message flow)};

\end{tikzpicture}

}
  \caption{Node graph of the autonomous indoor robot, with the key topics between the sensor, processing, navigation, and actuator nodes.}
  \label{fig:kap15-node-graph}
\end{figure}

\section{Granularity}
Granularity is an important decision. Too many nodes raise communication overhead. Too few nodes make the system hard to follow. As a starting rule: a clearly bounded function may become its own node. A node that handles the camera, the lidar, and navigation all at once is cut too coarse; a separate node for every single sensor-reading function is usually cut too fine.

\section{Architecture Decisions}
Key decisions include:

\begin{itemize}
    \item Which functions become their own nodes?
    \item Which data gets published as topics?
    \item Which processes are modeled as actions?
    \item Which parameters must be configurable?
    \item Which nodes are safety-relevant?
    \item Which errors must be reported centrally?
\end{itemize}

\section{Failure Cases}
An architecture should not show only the normal case. What happens if the camera fails? What happens when the battery is empty? What happens if no map is available?

\begin{praxisbox}
If \texttt{camera\_node} delivers no images, person detection must not simply keep reusing old results. \texttt{object\_detector\_node} should detect the missing data and report a matching status.
\end{praxisbox}

\section{Parameters}
Typical parameters include maximum speed, battery threshold, topic names, detection limits, or timeout values. Parameters should match the assumptions and requirements in the model.

\section{Lifecycle Nodes}
ROS2 supports lifecycle nodes: nodes with a defined state machine for controlled activation and deactivation. They have defined states such as unconfigured, inactive, active, and finalized. This is optional for beginners, but useful for robust systems: the state machine enables, for example, a deterministic startup in which a sensor node first loads its parameters (state \emph{configure}) before it actively delivers data (state \emph{active}) -- so navigation does not start with invalid or incomplete data. Figure~\ref{fig:kap15-lifecycle} shows the matching state machine.

\begin{figure}[H]
  \centering
  \adjustbox{max width=\linewidth,max totalheight=0.6\textheight}{% =====================================================================
% Figure: Lifecycle state machine of a ROS2 node
% Matches Chapter 15, section "Lifecycle Nodes"
%
% Included in the book, e.g. in chapters_en/15_ros2_architektur.tex:
%
%   \begin{figure}[H]
%     \centering
%     % =====================================================================
% Figure: Lifecycle state machine of a ROS2 node
% Matches Chapter 15, section "Lifecycle Nodes"
%
% Included in the book, e.g. in chapters_en/15_ros2_architektur.tex:
%
%   \begin{figure}[H]
%     \centering
%     % =====================================================================
% Figure: Lifecycle state machine of a ROS2 node
% Matches Chapter 15, section "Lifecycle Nodes"
%
% Included in the book, e.g. in chapters_en/15_ros2_architektur.tex:
%
%   \begin{figure}[H]
%     \centering
%     \input{figures_en/fig_kap15_lifecycle_zustandsautomat}
%     \caption{State machine of a ROS2 lifecycle node, from
%       \texttt{Unconfigured} through \texttt{Inactive} and \texttt{Active}
%       to the final state \texttt{Finalized}.}
%     \label{fig:kap15-lifecycle}
%   \end{figure}
%
% Required packages (already loaded in the book via main_en.tex): tikz
% Required TikZ libraries: positioning, arrows.meta, shadows, calc
% =====================================================================

\input{figures_en/fig_kap15_lifecycle_zustandsautomat.tikz}

%     \caption{State machine of a ROS2 lifecycle node, from
%       \texttt{Unconfigured} through \texttt{Inactive} and \texttt{Active}
%       to the final state \texttt{Finalized}.}
%     \label{fig:kap15-lifecycle}
%   \end{figure}
%
% Required packages (already loaded in the book via main_en.tex): tikz
% Required TikZ libraries: positioning, arrows.meta, shadows, calc
% =====================================================================

\begin{tikzpicture}[
    scale=0.78, transform shape,
    every node/.style={font=\small},
    state/.style={
        rectangle, rounded corners=2pt, draw=black!70, thick,
        minimum width=30mm, minimum height=10mm,
        align=center, fill=blue!10, drop shadow
    },
    finalstate/.style={state, fill=black!10, draw=black!60, double, double distance=1.2pt},
    initial/.style={circle, fill=black, minimum size=3mm, inner sep=0pt},
    trans/.style={-{Stealth[length=2.2mm]}, thick, black!70},
    lbl/.style={font=\small, fill=white, inner sep=1pt}
]

% --- States ---
\node[initial]    (init)   at (-2.6, 0)   {};
\node[state]      (unconf) at (0, 0)      {\texttt{Unconfigured}};
\node[state]      (inact)  at (5.5, 0)    {\texttt{Inactive}};
\node[state]      (act)    at (11, 0)     {\texttt{Active}};
\node[finalstate] (fin)    at (5.5, -5)   {\texttt{Finalized}};

% --- entry ---
\draw[trans] (init) -- (unconf);

% --- Main cycle ---
\draw[trans] (unconf.east) -- node[lbl, above]{configure} (inact.west);
\draw[trans] ([yshift=2mm]inact.east) -- node[lbl, above]{activate} ([yshift=2mm]act.west);
\draw[trans] ([yshift=-2mm]act.west) -- node[lbl, below]{deactivate} ([yshift=-2mm]inact.east);
\draw[trans] (inact.north) .. controls ++(0,1.6) and ++(0,1.6) .. node[lbl, above, pos=0.5]{cleanup} (unconf.north);

% --- shutdown transitions to Finalized ---
\draw[trans] (unconf.south) -- ++(0,-2.2) -| node[lbl, pos=0.55, above]{shutdown} (fin.north west);
\draw[trans] (inact.south) -- node[lbl, right]{shutdown} (fin.north);

% --- Legend ---
\node[state, minimum width=6mm, minimum height=4mm, font=\small] (leg1) at (0,-8) {};
\node[right=1mm of leg1, font=\small, align=left] (lab1){Normal state};
\node[finalstate, minimum width=6mm, minimum height=4mm, font=\small, right=6mm of lab1] (leg2) {};
\node[right=1mm of leg2, font=\small, align=left] (lab2){Final state (\texttt{Finalized})};
\node[initial, right=6mm of lab2] (leg3) {};
\node[right=1mm of leg3, font=\small, align=left] (lab3){Start state};

\end{tikzpicture}


%     \caption{State machine of a ROS2 lifecycle node, from
%       \texttt{Unconfigured} through \texttt{Inactive} and \texttt{Active}
%       to the final state \texttt{Finalized}.}
%     \label{fig:kap15-lifecycle}
%   \end{figure}
%
% Required packages (already loaded in the book via main_en.tex): tikz
% Required TikZ libraries: positioning, arrows.meta, shadows, calc
% =====================================================================

\begin{tikzpicture}[
    scale=0.78, transform shape,
    every node/.style={font=\small},
    state/.style={
        rectangle, rounded corners=2pt, draw=black!70, thick,
        minimum width=30mm, minimum height=10mm,
        align=center, fill=blue!10, drop shadow
    },
    finalstate/.style={state, fill=black!10, draw=black!60, double, double distance=1.2pt},
    initial/.style={circle, fill=black, minimum size=3mm, inner sep=0pt},
    trans/.style={-{Stealth[length=2.2mm]}, thick, black!70},
    lbl/.style={font=\small, fill=white, inner sep=1pt}
]

% --- States ---
\node[initial]    (init)   at (-2.6, 0)   {};
\node[state]      (unconf) at (0, 0)      {\texttt{Unconfigured}};
\node[state]      (inact)  at (5.5, 0)    {\texttt{Inactive}};
\node[state]      (act)    at (11, 0)     {\texttt{Active}};
\node[finalstate] (fin)    at (5.5, -5)   {\texttt{Finalized}};

% --- entry ---
\draw[trans] (init) -- (unconf);

% --- Main cycle ---
\draw[trans] (unconf.east) -- node[lbl, above]{configure} (inact.west);
\draw[trans] ([yshift=2mm]inact.east) -- node[lbl, above]{activate} ([yshift=2mm]act.west);
\draw[trans] ([yshift=-2mm]act.west) -- node[lbl, below]{deactivate} ([yshift=-2mm]inact.east);
\draw[trans] (inact.north) .. controls ++(0,1.6) and ++(0,1.6) .. node[lbl, above, pos=0.5]{cleanup} (unconf.north);

% --- shutdown transitions to Finalized ---
\draw[trans] (unconf.south) -- ++(0,-2.2) -| node[lbl, pos=0.55, above]{shutdown} (fin.north west);
\draw[trans] (inact.south) -- node[lbl, right]{shutdown} (fin.north);

% --- Legend ---
\node[state, minimum width=6mm, minimum height=4mm, font=\small] (leg1) at (0,-8) {};
\node[right=1mm of leg1, font=\small, align=left] (lab1){Normal state};
\node[finalstate, minimum width=6mm, minimum height=4mm, font=\small, right=6mm of lab1] (leg2) {};
\node[right=1mm of leg2, font=\small, align=left] (lab2){Final state (\texttt{Finalized})};
\node[initial, right=6mm of lab2] (leg3) {};
\node[right=1mm of leg3, font=\small, align=left] (lab3){Start state};

\end{tikzpicture}

}
  \caption{State machine of a ROS2 lifecycle node, from \texttt{Unconfigured} through \texttt{Inactive} and \texttt{Active} to the final state \texttt{Finalized}.}
  \label{fig:kap15-lifecycle}
\end{figure}

\section{Architecture Review}
Before you write any code, briefly review the ROS2 architecture. A review can follow simple questions:

\begin{itemize}
    \item Does every node have a clear responsibility?
    \item Do topics have a clear sender and receiver?
    \item Are safety-relevant data flows identifiable?
    \item Are failure cases described?
    \item Are parameters provided for values that can change?
    \item Is there a link to requirements and tests?
\end{itemize}

The tool \texttt{rqt\_graph} can help you check the designed node structure visually later, in the running system.

For the indoor robot, it is especially important that navigation, obstacle detection, and motion control are not coupled without control. A wrong motion command can have an immediate physical effect. The data flow to \texttt{/cmd\_vel} should therefore stay clearly traceable.

\section{Summary}
A good ROS2 architecture has clear responsibilities, understandable communication, and defined error responses. You should be able to derive it from the SysML model, and later check it through tests.

\section{Review Questions}
\begin{enumerate}
    \item What does a ROS2 architecture describe?
    \item Why does node granularity matter?
    \item Which nodes does the indoor robot need at minimum?
    \item Which topics matter for navigation?
    \item Why must an architecture consider failure cases?
    \item What are parameters?
    \item What are lifecycle nodes?
\end{enumerate}

\section{Exercises}
\begin{uebungbox}
\textbf{Exercise 1: Sketch an Architecture}\\
Create a ROS2 architecture sketch with at least six nodes and eight topics.
\end{uebungbox}

\begin{uebungbox}
\textbf{Exercise 2: Analyze a Failure Case}\\
Describe which nodes are affected if the lidar fails. Which status message should the system generate?
\end{uebungbox}

\section{Further Reading}
The ROS2 documentation \cite{ros2docs} describes communication mechanisms, lifecycle nodes, launch files, and tools. For architecture decisions, the SysML model from the previous chapters is also helpful.
